\documentclass[11pt]{amsart}
\usepackage{../cqts-formalization-series}
\newcommand{\CH}{\mathbb{C}_{\mathrm H}}
\newcommand{\CHunits}{\mathbb{C}_{\mathrm H}^{\times}}
\newcommand{\BG}{\mathbf{B}}
\newcommand{\Loop}{\mathop{\Omega}}
\newcommand{\Type}{\mathsf{Type}}
\title[Constructive Complex Units and Their Torsor Delooping]{Constructive Complex Units and Their\\[0.35em]Torsor Delooping in Cubical Agda\\[0.35em]\large\textnormal{A Mathematical Companion to a CQTS Formalization}}
\author{Hongwei (Alex) Wang\\[3pt]
\small New York University Abu Dhabi\\
\small \texttt{hw4375@nyu.edu}}
\dedicatory{\normalfont\normalsize 23 November 2026}

\begin{document}
\begin{abstract}
We document a Cubical Agda construction of constructive complex numbers from
Jackson Brough's HoTT real numbers, their algebraic multiplicative units, and
a torsor-moduli delooping.  The local code proves that
$\CH=\mathbb{R}_{\mathrm H}\times\mathbb{R}_{\mathrm H}$ is a commutative
ring, packages its algebraic units as a group $\CHunits$, and specializes an
existing univalent torsor construction to a pointed type $\BG\CHunits$.
The checked endpoint is an equivalence of underlying types
$\Loop(\BG\CHunits)\simeq\CHunits$ at the regular-torsor basepoint.  This
article supplies mathematical proofs and declaration-level provenance while
keeping the boundary explicit: it proves neither a coordinate-level
apartness criterion for units, nor a packaged group isomorphism on loops,
nor dependent computational transport or quantum gates.
\end{abstract}
\maketitle
\pagestyle{plain}

\section{Introduction and mathematical motivation}
\label{sec:introduction}

The CQTS project investigates whether suitably homotopical topological quantum
transformations can acquire an executable realization as dependent-type
transport in Cubical Agda.  The motivating perspective is the programme of
Myers--Sati--Schreiber \cite{mss}.  Its later stages require precise higher
types, dependent families over them, loops, and a computation analysis of
transport.  None of those later constructions follows merely from naming a
familiar group or classifying space.

This paper records a completed smaller layer needed to formulate one possible
such direction.  Jackson Brough's constructive HoTT real numbers provide a
checked real-number object $\mathbb{R}_{\mathrm H}$.  From it the local
development constructs
\begin{equation}
  \CH\coloneqq\mathbb{R}_{\mathrm H}\times\mathbb{R}_{\mathrm H},
\end{equation}
proves that it is a commutative ring, takes its algebraic units $\CHunits$,
and applies the existing generic theory of left torsors.  The result is the
pointed torsor-moduli type
\begin{equation}
 (\BG\CHunits,b_{\CHunits})\quad\text{with}\quad
 \Loop(\BG\CHunits,b_{\CHunits})\simeq\CHunits .
 \label{eq:main-result}
\end{equation}

The notation $\BG\CHunits$ is a name for this torsor-moduli construction, not
a claim that a separate conventional geometric classifying-space model has
been compared with it.  Equation~\eqref{eq:main-result} is the exact
type-equivalence endpoint of the formalization.  It is neither a newly
packaged group isomorphism nor a construction of a quantum gate.

The argument is layered.  Section~\ref{sec:complex} gives the constructive
complex ring; Section~\ref{sec:units} its algebraic unit group; and
Sections~\ref{sec:torsors}--\ref{sec:loops} specialize the existing
univalent torsor delooping and trace its loop computation.  The provenance
map and reproducibility record are included in the appendices.

\section{Constructive complex numbers}
\label{sec:complex}

\subsection{Carrier and operations}
Let $R=\mathbb{R}_{\mathrm H}$.  The formal carrier is the product
$\CH=R\times R$.  Its operations are
\begin{align}
0_{\CH}&=(0,0), & 1_{\CH}&=(1,0),\\
(a,b)+(c,d)&=(a+c,b+d), & -(a,b)&=(-a,-b),\\
(a,b)(c,d)&=(ac-bd,ad+bc),\\
i&=(0,1), & \overline{(a,b)}&=(a,-b).
\end{align}
Thus the imaginary unit and conjugation are checked definitions, although no
separate theorem about $i^2$ or conjugation multiplicativity is claimed by
these modules.

\formal{\texttt{src/ConstructiveComplex/Complex.agda}}
{\texttt{Complex}; \texttt{complexZero}; \texttt{complexOne}; \texttt{complexAdd}; \texttt{complexNeg}; \texttt{complexMul}; \texttt{imaginaryUnit}; \texttt{complexConjugate}}
{These are definitions, with the displayed coordinate formula for multiplication.}

The only HoTT Reals declarations used directly are the carrier
\agda{HoTTReals.Data.Real.Base.ℝ}, the set proof
\agda{RealProperties.ℝ-isSet}, and the commutative-ring package
\agda{RealMultiplication.ℝCommutativeRing}.  The local development does not
reproduce Brough's construction of real numbers.

\upstream{\texttt{HoTTReals/Data/Real/Base.agda}}{\texttt{ℝ}}
\upstream{\texttt{HoTTReals/Data/Real/Properties.agda}}{\texttt{ℝ-isSet}}
\upstream{\texttt{HoTTReals/Data/Real/Algebra/Multiplication.agda}}{\texttt{ℝCommutativeRing}}

\subsection{Sethood and pair paths}
Since $R$ is a set, its product is a set by the imported closure theorem
\agda{isSet×}.  Equalities of complex numbers are assembled from equalities
of real coordinates with \agda{ΣPathP}:
\begin{equation}
 a=b\ \text{and}\ c=d\quad\Longrightarrow\quad(a,c)=(b,d).
\end{equation}
\formal{\texttt{src/ConstructiveComplex/Complex.agda}}
{\texttt{complexIsSet}; \texttt{pairPath}}
{\texttt{complexIsSet} uses \texttt{isSet×}; the private helper \texttt{pairPath} is \texttt{ΣPathP}.}

\subsection{The ring proof}
\begin{proposition}
The displayed operations make $\CH$ a commutative ring.
\end{proposition}
\begin{proof}
All eight algebraic laws reduce to polynomial identities in the commutative
ring $R$.  For example, associativity compares
\begin{align*}
(a,b)((c,d)(e,f))&=\bigl(a(ce-df)-b(cf+de),\ a(cf+de)+b(ce-df)\bigr),\\
((a,b)(c,d))(e,f)&=\bigl((ac-bd)e-(ad+bc)f,\ (ac-bd)f+(ad+bc)e\bigr).
\end{align*}
The two coordinates agree by associativity, commutativity, distributivity,
and additive-inverse laws in $R$; \agda{ΣPathP} then gives equality in
$\CH$.  Addition, the unit and inverse laws, distributivity, and
commutativity of multiplication are proved by the same coordinatewise
method.  The sethood proof and these eight fields are precisely the data
accepted by Cubical's commutative-ring constructor.
\end{proof}

The formal proof discharges each coordinate by the verified commutative-ring
solver instantiated at Brough's ring, rather than treating the displayed
calculation as an unchecked external argument:
\begin{verbatim}
complexMulAssociative (a , b) (c , d) (e , f) =
  pairPath (solve! RealMultiplication.ℝCommutativeRing)
           (solve! RealMultiplication.ℝCommutativeRing)
\end{verbatim}
\formal{\texttt{src/ConstructiveComplex/Complex.agda}}
{\texttt{complexAddAssociative}; \texttt{complexAddUnitRight}; \texttt{complexAddInverseRight}; \texttt{complexAddCommutative}; \texttt{complexMulAssociative}; \texttt{complexMulUnitRight}; \texttt{complexMulDistributesOverAddLeft}; \texttt{complexMulCommutative}; \texttt{complexIsCommRing}; \texttt{ComplexCommRing}}
{The eight local laws are combined by \texttt{makeIsCommRing}; \texttt{ComplexCommRing} is the resulting level-zero package.}

\section{The multiplicative group of complex units}
\label{sec:units}

\subsection{Algebraic units as a subtype}
For a commutative ring $A$, Cubical's $A^{\times}$ is the predicate of
elements equipped with multiplicative-inverse evidence.  The local definition
is the dependent sum
\begin{equation}
 \CHunits\coloneqq\sum_{z:\CH}(z\in\CH^{\times}).
 \label{eq:units-sigma}
\end{equation}
An element is a complex number together with algebraic evidence that it is a
unit.  This uses neither a quotient nor a classical coordinate test for
nonzero values.
\formal{\texttt{src/ConstructiveComplex/Units.agda}}
{\texttt{ComplexUnits}}
{The definition is \texttt{Σ Complex U.Rˣ}, where \texttt{U = Units ComplexCommRing}.}

Membership in the unit predicate is propositional by \agda{∈-isProp}.
Therefore two unit records with equal first projections are equal, by
\agda{Σ≡Prop}; and since $\CH$ is a set, the sigma type is a set by
\agda{isSetΣSndProp}.
\upstream{\texttt{Cubical/Foundations/Powerset.agda}}{\texttt{∈-isProp}}
\upstream{\texttt{Cubical/Foundations/HLevels.agda}}{\texttt{isSetΣSndProp}; \texttt{Σ≡Prop}}
\formal{\texttt{src/ConstructiveComplex/Units.agda}}
{\texttt{complexUnitsIsSet}; \texttt{unitPath}}
{The private \texttt{unitPath} lifts equality of underlying complex numbers to equality of units.}

\subsection{Group operations and proof}
Writing a unit as $(z,p)$, the source defines
\begin{align}
1_{\CHunits}&=(1,\text{unit evidence}),\\
(z,p)(w,q)&=(zw,\text{closure evidence}),\\
(z,p)^{-1}&=(z^{-1},\text{closure evidence}).
\end{align}
The generic \agda{Units} API supplies the evidence
\agda{RˣContainsOne}, \agda{RˣMultClosed}, and \agda{RˣInvClosed}; its
\agda{·-rinv} establishes $zz^{-1}=1$ for a certified unit.
\upstream{\texttt{Cubical/Algebra/CommRing/Properties.agda}}{\texttt{Units.RˣContainsOne}; \texttt{Units.RˣMultClosed}; \texttt{Units.RˣInvClosed}; \texttt{Units.·-rinv}}
\formal{\texttt{src/ConstructiveComplex/Units.agda}}
{\texttt{complexUnitOne}; \texttt{complexUnitMul}; \texttt{complexUnitInv}}
{These definitions instantiate the generic unit API at \texttt{ComplexCommRing}.}

\begin{proposition}
$\CHunits$ is a group under the displayed multiplication.
\end{proposition}
\begin{proof}
For $u=(z,p)$, $v=(w,q)$, and $r=(t,s)$, associativity becomes
$z(wt)=(zw)t$ after \agda{Σ≡Prop}, hence is
\agda{complexMulAssociative}.  The right unit law similarly becomes
\agda{complexMulUnitRight}, and the generic inverse theorem gives the right
inverse.  Cubical's \agda{makeGroup-right} packages sethood and these three
right-handed laws into a group; its record supplies the derived left unit and
left inverse laws.
\end{proof}
\formal{\texttt{src/ConstructiveComplex/Units.agda}}
{\texttt{complexUnitMulAssociative}; \texttt{complexUnitMulUnitRight}; \texttt{complexUnitMulInverseRight}; \texttt{ComplexUnitsGroup}; \texttt{complexUnitMulUnitLeft}; \texttt{complexUnitMulInverseLeft}}
{The first three are inputs to \texttt{makeGroup-right}; the latter two are derived from the resulting \texttt{Group} record.}

\begin{warning}[Boundary of the unit construction]
No declaration proves an equivalence between algebraic unit membership and a
coordinate-level apartness or nonzero condition on $(a,b)$.  In particular,
the formalization does not claim a constructive coordinate criterion for
membership in $\CHunits$.
\end{warning}

\section{Regular torsors and delooping}
\label{sec:torsors}

\subsection{Left torsors and the regular example}
For a set-level group $G$, a left torsor is a left $G$-action on a type $X$
which is merely inhabited and has contractible transporter types
\begin{equation}
  \sum_{g:G}(g\cdot x=y)
\end{equation}
for all $x,y:X$.  The second condition is the type-theoretic free-and-
transitive condition; the truncation in the first condition is important,
since a general torsor has no selected origin.
\formal{\texttt{src/TorsorDelooping/LeftTorsor.agda}}
{\texttt{IsLeftTorsor}; \texttt{LeftTorsor}}
{The fields are \agda{inhabited : ∥ X ∥₁} and the contractible-transporter field \texttt{regular}.}

The regular left torsor is $G$ itself under left multiplication.  It is
inhabited by $1$, and its transporter is contractible through the imported
right-translation equivalence.  The use of right translations is compatible
with the left action: it proves the needed regularity without changing the
chosen action convention.
\formal{\texttt{src/TorsorDelooping/Regular.agda}}
{\texttt{regularLeftAction}; \texttt{regularLeftTorsor}}
{Regularity uses \texttt{·GroupAutomorphismR}; the inhabitant is the group unit.}

\subsection{The univalent moduli type}
The generic construction first records a left action, then adds the
propositional fields asserting sethood, the action laws, mere inhabitation,
and contractible transporters.  The Structure Identity Principle therefore
identifies paths of structured torsors with equivariant equivalences.  It
constructs a moduli type $G\mathrm{Tors}_{L,\ell_X}$ and an equivalence
\begin{equation}
 \operatorname{TorsorEquiv}(T,U)\simeq(T=U).
 \label{eq:sip-torsors}
\end{equation}
\formal{\texttt{src/TorsorDelooping/Moduli.agda}}
{\texttt{LeftTorsorModuli}; \texttt{TorsorEquiv}; \texttt{torsorPathEquiv}; \texttt{fromLeftTorsor}}
{\texttt{torsorPathEquiv} is the SIP equivalence; \texttt{fromLeftTorsor} embeds the regular-torsor record into the moduli type.}

\subsection{The constructive complex basepoint}
Set $G=\CHunits$.  The local units module supplies the regular torsor
\agda{complexUnitsRegularLeftTorsor}; its image in the moduli type defines
\begin{align}
 \BG\CHunits&\coloneqq\CHunits\mathrm{Tors}_{L,0},\\
 b_{\CHunits}&\coloneqq\text{the regular left $\CHunits$-torsor}.
\end{align}
The source records \agda{BComplexUnits : Type₁}; that universe level is
inherited from the generic structured-type construction.
\formal{\texttt{src/ConstructiveComplex/Units.agda}}
{\texttt{complexUnitsRegularLeftTorsor}}
{This is the regular torsor of \texttt{ComplexUnitsGroup}.}
\formal{\texttt{src/ConstructiveComplex/Delooping.agda}}
{\texttt{BComplexUnits}; \texttt{bComplexUnits}}
{Both are direct specializations of \texttt{LeftTorsorModuli} and \texttt{torsorBase}.}

\section{Loop-space recovery}
\label{sec:loops}

\subsection{Loops as equivariant automorphisms}
At the regular basepoint, the inverse of \eqref{eq:sip-torsors} gives
\begin{equation}
 \Loop(G\mathrm{Tors}_{L,\ell},b_G)\simeq\operatorname{Aut}_G(G),
 \label{eq:loops-aut}
\end{equation}
where $G$ has the regular left action.  In the source this is
\agda{torsorLoopEquivAut}, obtained by applying \agda{invEquiv} to
\agda{torsorPathEquiv}.
\formal{\texttt{src/TorsorDelooping/Delooping.agda}}
{\texttt{torsorLoopEquivAut}}
{The equivalence is from loops to structured equivariant self-equivalences of the regular torsor.}

\subsection{Evaluation and right translations}
Evaluation at the unit is the map $E(f)=f(1)$ from equivariant automorphisms
to $G$.  Its inverse sends $h:G$ to the right translation
\begin{equation}
 R_h(x)=xh.
 \label{eq:right-translation}
\end{equation}
It is equivariant for the left action because $R_h(gx)=(gx)h=g(xh)$.  It is
invertible through $R_{h^{-1}}$.  Conversely, equivariance forces
\begin{equation}
 f(x)=f(x1)=xf(1)=R_{f(1)}(x).
 \label{eq:equivariant-determined}
\end{equation}
Function extensionality proves equality of the underlying equivalences, and
\agda{Σ≡Prop} lifts it through the propositional equivariance field.  Thus
evaluation and right translation form an isomorphism.
\formal{\texttt{src/TorsorDelooping/Delooping.agda}}
{\texttt{regularAutToGroup}; \texttt{rightRegularAut}; \texttt{regularAutUnderlying}; \texttt{regularAutRet}; \texttt{regularAutIsoGroup}}
{These declarations implement evaluation, right translation, extensional reconstruction, and the resulting \texttt{Iso}.}

Composing \eqref{eq:loops-aut} with this isomorphism yields
\begin{equation}
 \Loop(G\mathrm{Tors}_{L,\ell},b_G)\simeq G.
 \label{eq:generic-loop}
\end{equation}
\formal{\texttt{src/TorsorDelooping/Delooping.agda}}
{\texttt{torsorLoopEquivGroup}}
{Its implementation is \texttt{compEquiv torsorLoopEquivAut (isoToEquiv regularAutIsoGroup)}.}

\subsection{Order conventions and the formal boundary}
There is no hidden opposite group in this calculation.  If composition means
first $R_g$ and then $R_h$, then
\begin{equation}
 R_h(R_g(x))=(xg)h=x(gh)=R_{gh}(x).
\end{equation}
This is the correct order for right translations acting as equivariant maps
of a left torsor.  The generic CQTS development has a separate regular-loop
coherence used for its $\mathsf{EM}_1$ construction, but
\agda{torsorLoopEquivGroup} is not packaged as a group-isomorphism theorem.

\begin{warning}[No group-level loop-recovery claim]
Loop concatenation gives the source of \eqref{eq:generic-loop} a group
structure, but the formal endpoint supplies only an equivalence of underlying
types.  No formal theorem here asserts preservation of multiplication, unit,
or inverse.
\end{warning}

For $G=\CHunits$, \eqref{eq:generic-loop} is the local declaration
\begin{equation}
 \Loop(\BG\CHunits,b_{\CHunits})\simeq\CHunits.
\end{equation}
\formal{\texttt{src/ConstructiveComplex/Delooping.agda}}
{\texttt{complexUnitsLoopEquiv}}
{This is definitionally \texttt{torsorLoopEquivGroup ComplexUnitsGroup}; its codomain is \texttt{ComplexUnits}.}

\section{Formalization and proof provenance}
\label{sec:provenance}

The construction has four distinct sources of proof content.
\begin{enumerate}[label=\textbf{\arabic*.}]
\item \textbf{HoTT Reals.} Brough's \texttt{hott-reals} thesis-branch revision
supplies $\mathbb{R}_{\mathrm H}$, its sethood, and its commutative-ring
package.  No upstream real-number source is copied into this repository.
\item \textbf{Cubical Agda.} The pinned library supplies sigma paths, h-level
closure, the commutative-ring solver, the algebraic unit API, group packaging,
univalence, SIP, and right-translation equivalences.
\item \textbf{Existing CQTS torsor formalization.} The modules
\texttt{TorsorDelooping.Regular}, \texttt{Moduli}, and \texttt{Delooping}
prove the generic regular torsor, univalent moduli, and loop recovery used
here.
\item \textbf{New ConstructiveComplex layer.} The three local modules define
the pair ring, its units group, and their specialization of the generic
delooping.
\end{enumerate}

The following is the declaration-level correspondence; it is provenance, not
an additional collection of claims.

\begingroup\small
\begin{longtable}{@{}p{0.28\textwidth}p{0.29\textwidth}p{0.33\textwidth}@{}}
\toprule
Mathematical statement & Agda declaration and source & Direct proof dependencies \\
\midrule
\endhead
$\CH=R\times R$ with complex operations &
\agda{Complex}, \agda{complexAdd}, \agda{complexMul};
\texttt{Complex.agda} &
Brough's \agda{ℝ}; definitions of product operations. \\

$\CH$ is a commutative ring &
\agda{complexIsCommRing}, \agda{ComplexCommRing};
\texttt{Complex.agda} &
\agda{ℝ-isSet}, \agda{ℝCommutativeRing}, \agda{isSet×},
\agda{ΣPathP}, \agda{solve!}, \agda{makeIsCommRing}. \\

Algebraic complex units form a set &
\agda{complexUnitsIsSet}; \texttt{Units.agda} &
\agda{complexIsSet}, \agda{∈-isProp}, \agda{isSetΣSndProp}. \\

$\CHunits$ is a group &
\agda{ComplexUnitsGroup}; \texttt{Units.agda} &
Generic \agda{Units} closure/inverse facts, local complex laws, and
\agda{makeGroup-right}. \\

Regular left $\CHunits$-torsor &
\agda{complexUnitsRegularLeftTorsor}; \texttt{Units.agda} &
\agda{regularLeftTorsor} from \texttt{TorsorDelooping.Regular}. \\

Torsor moduli and paths &
\agda{LeftTorsorModuli}, \agda{torsorPathEquiv};
\texttt{Moduli.agda} &
Cubical structured types, propositional axioms, and SIP. \\

Generic loop recovery &
\agda{torsorBase}, \agda{torsorLoopEquivGroup};
\texttt{Delooping.agda} &
Regular torsor, SIP, \agda{·GroupAutomorphismR}, evaluation, and right
translation. \\

$\BG\CHunits$, basepoint, loop equivalence &
\agda{BComplexUnits}, \agda{bComplexUnits},
\agda{complexUnitsLoopEquiv}; \texttt{ConstructiveComplex/Delooping.agda} &
Direct specialization at \agda{ComplexUnitsGroup}. \\
\bottomrule
\end{longtable}
\endgroup

The recorded environment is Agda~2.8.0--3d04bac, with Cubical revision
\begin{center}
\texttt{87af036ebd1c87956cf08fea44e62c7cf76392ca}
\end{center}
and Brough's \texttt{hott-reals} revision
\begin{center}
\texttt{9fcc142b86d0c1e2bd6d06f86f7617200ad74857}.
\end{center}
Those revisions identify the interfaces cited above; the paper neither vendors
nor modifies either upstream project.

\section{Scope, limitations, and future work}
\label{sec:scope}

\subsection{What is established}
The checked endpoint is concrete: a constructive complex commutative ring
$\CH$, an algebraic group of units $\CHunits$, its regular left torsor, the
torsor-moduli type $\BG\CHunits$, its canonical point, and the underlying
type equivalence of Equation~\eqref{eq:main-result}.  The generic theory also
establishes connectedness and one-truncation for the torsor moduli, though the
constructive-complex module uses only the basepoint and loop equivalence.

\subsection{Explicit non-results}
The following statements are outside the current formal endpoint.
\begin{enumerate}[label=\arabic*.]
\item No coordinate-level apartness/nonzero characterization of algebraic
complex units is proved.
\item No group-level loop recovery theorem is packaged for
\agda{complexUnitsLoopEquiv}.
\item No general proof identifies complex-unit multiplication with loop
concatenation for that displayed equivalence.
\item No dependent-family transport computation or monodromy calculation is
constructed.
\item No quantum-gate realization is implemented.
\item Myers--Sati--Schreiber Definition~6.7 is not formalized.
\end{enumerate}

These are not cosmetic omissions.  A group isomorphism requires a stated loop
group structure and proofs preserving product, unit, and inverse.  A quantum
transformation requires a specified dependent family, a loop or higher action
on it, and an analysis of the corresponding transport.  Neither structure
follows merely from the type equivalence already proved.

\subsection{Next mathematical stages}
A later development may choose a constructive apartness relation for complex
numbers and prove its relation to algebraic units.  It may then formulate a
group-level refinement of loop recovery with explicit orientation conventions.
Only after specifying a dependent family over a candidate delooping could the
project study transport, its computational behaviour, and its possible
relationship with the topological-quantum motivation of \cite{mss}.  Those
would be new formalization tasks rather than consequences of this paper.

\begin{thebibliography}{9}
\bibitem{hottbook}
The Univalent Foundations Program.
\newblock \emph{Homotopy Type Theory: Univalent Foundations of Mathematics}.
\newblock Institute for Advanced Study, 2013.
\newblock \url{https://homotopytypetheory.org/book/}.

\bibitem{brough}
Jackson Brough.
\newblock \emph{hott-reals}, thesis branch, revision
\texttt{9fcc142b86d0c1e2bd6d06f86f7617200ad74857}.
\newblock \url{https://github.com/utahplt/hott-reals}.

\bibitem{cubical}
Cubical Agda contributors.
\newblock \emph{Cubical Agda standard library}, v0.9-compatible line.
\newblock Pinned project checkout:
\texttt{87af036ebd1c87956cf08fea44e62c7cf76392ca}.
\newblock \url{https://github.com/agda/cubical}.

\bibitem{agda}
The Agda Team.
\newblock \emph{Agda 2.8.0-3d04bac}.
\newblock \url{https://agda.readthedocs.io/}.

\bibitem{mss}
David Jaz Myers, Hisham Sati, and Urs Schreiber.
\newblock \emph{Topological Quantum Gates in Homotopy Type Theory}.
\newblock \emph{Communications in Mathematical Physics} \textbf{405},
Article 172 (2024).
\newblock \url{https://doi.org/10.1007/s00220-024-05020-8}.
\end{thebibliography}

\appendix
\section{Statement-to-declaration index}
\label{app:index}

This appendix gives repository-relative identifiers for the claims in the
main text.  Upstream denotes the pinned Cubical or separate HoTT Reals
checkout; generic CQTS denotes the pre-existing torsor development.

\begingroup\small\raggedright
\begin{description}[font=\normalfont\bfseries]
\item[HoTT Reals (upstream)]
\texttt{HoTTReals/Data/Real/Base.agda}: \agda{ℝ}.
\texttt{HoTTReals/Data/Real/Properties.agda}: \agda{ℝ-isSet}.
\texttt{HoTTReals/Data/Real/Algebra/Multiplication.agda}:
\agda{ℝCommutativeRing}.

\item[Constructive complex ring (new local)]
\texttt{src/ConstructiveComplex/Complex.agda}: \agda{Complex},
\agda{complexZero}, \agda{complexOne}, \agda{complexAdd},
\agda{complexNeg}, \agda{complexMul}, \agda{imaginaryUnit},
\agda{complexConjugate}, \agda{complexIsSet}, \agda{complexIsCommRing},
\agda{ComplexCommRing}.

\item[Algebraic units (new local)]
\texttt{src/ConstructiveComplex/Units.agda}: \agda{ComplexUnits},
\agda{complexUnitOne}, \agda{complexUnitMul}, \agda{complexUnitInv},
\agda{complexUnitsIsSet}, \agda{ComplexUnitsGroup},
\agda{complexUnitsRegularLeftTorsor}.

\item[Generic torsor construction (CQTS)]
\texttt{src/TorsorDelooping/Regular.agda}: \agda{regularLeftAction},
\agda{regularLeftTorsor}.  \texttt{src/TorsorDelooping/Moduli.agda}:
\agda{LeftTorsorModuli}, \agda{TorsorEquiv}, \agda{torsorPathEquiv},
\agda{fromLeftTorsor}.

\item[Generic loop recovery (CQTS)]
\texttt{src/TorsorDelooping/Delooping.agda}: \agda{torsorBase},
\agda{torsorLoopEquivAut}, \agda{regularAutToGroup},
\agda{rightRegularAut}, \agda{regularAutIsoGroup},
\agda{torsorLoopEquivGroup}.

\item[Constructive-complex delooping (new local)]
\texttt{src/ConstructiveComplex/Delooping.agda}: \agda{BComplexUnits},
\agda{bComplexUnits}, \agda{complexUnitsLoopEquiv}.

\item[Key Cubical interfaces (upstream)]
\texttt{Cubical/Algebra/CommRing/Base.agda}: \agda{makeIsCommRing}.
\texttt{Cubical/Tactics/CommRingSolver.agda}: \agda{solve!}.
\texttt{Cubical/Algebra/CommRing/Properties.agda}: \agda{Units}.
\texttt{Cubical/Algebra/Group/Base.agda}: \agda{makeGroup-right}.
\texttt{Cubical/Algebra/Group/Properties.agda}: \agda{·GroupAutomorphismR}.
\texttt{Cubical/Foundations/SIP.agda}: \agda{SIP}.
\end{description}
\endgroup

The mathematical explanations in the main text explain checked constructions;
they do not assert that every sentence is a separately named theorem.

\section{Reproducibility and publication record}
\label{app:reproducibility}

\subsection{Formal environment}
The recorded environment is Agda \texttt{2.8.0-3d04bac}, the project-local
Cubical checkout at
\begin{center}
\texttt{87af036ebd1c87956cf08fea44e62c7cf76392ca},
\end{center}
and Brough's external \texttt{hott-reals} thesis branch at
\begin{center}
\texttt{9fcc142b86d0c1e2bd6d06f86f7617200ad74857}.
\end{center}
The project descriptor records the dependencies \texttt{cubical-0.9} and
\texttt{hott-reals}, with
\texttt{--safe --cubical --no-import-sorts --guardedness}.

After registering the three library descriptors described in
\texttt{src/ConstructiveComplex/README.md}, replay the checks sequentially:
{\small\begin{verbatim}
agda --safe --cubical --no-import-sorts --guardedness \
  --library=urs-group-presentation src/ConstructiveComplex/Complex.agda
agda --safe --cubical --no-import-sorts --guardedness \
  --library=urs-group-presentation src/ConstructiveComplex/Units.agda
agda --safe --cubical --no-import-sorts --guardedness \
  --library=urs-group-presentation src/ConstructiveComplex/Delooping.agda
\end{verbatim}}
The commands are intentionally sequential.  This paper does not run an Agda
check as part of its LaTeX build.

\subsection{Paper build}
From this directory, build the companion with:
{\small\begin{verbatim}
latexmk -xelatex -interaction=nonstopmode -halt-on-error main.tex
\end{verbatim}}
The PDF is a derived documentation artifact.  It neither changes Agda source
nor rebuilds dependency interfaces.

\subsection{Licensing and attribution}
The repository's \texttt{LICENSE} is MIT and applies to this new local
documentation.  Cubical Agda and \texttt{hott-reals} remain external
dependencies with their own repositories and provenance.  The paper cites
their roles and exact revisions but vendors neither upstream source.

\end{document}
